% 原创TikZ重绘；层次布局参考Bishop与Bishop（2024）第12章。
\begin{tikzpicture}[x=1mm,y=1mm,font=\figurefont\footnotesize,
 every node/.style={text=NNInk},
 stage/.style={nn operator,minimum height=9mm,inner sep=2pt},
 norm/.style={stage,minimum height=7mm},
 wire/.style={nn flow,draw=NNWire,rounded corners=1.5mm},
 aux/.style={draw=NNLine,line width=.85pt,densely dashed},
 sum/.style={nn pointwise,minimum size=5mm}]

\node[] (stoken0) at (17,0) {猫};
\node[inner sep=0pt,minimum width=12mm,minimum height=3mm] (sembed0) at (17,13) {};
\filldraw[draw=NNInput,fill=NNInput!18,line width=.5pt] (11.0,11.5) rectangle ++(3.0,3);
\filldraw[draw=NNInput,fill=NNInput!18,line width=.5pt] (14.0,11.5) rectangle ++(3.0,3);
\filldraw[draw=NNInput,fill=NNInput!18,line width=.5pt] (17.0,11.5) rectangle ++(3.0,3);
\filldraw[draw=NNInput,fill=NNInput!18,line width=.5pt] (20.0,11.5) rectangle ++(3.0,3);
\draw[wire] (stoken0.north)--(sembed0.south);
\draw[wire] (sembed0.north)--(17,51);
\node[] (stoken1) at (35,0) {$\langle S_1\rangle$};
\node[inner sep=0pt,minimum width=12mm,minimum height=3mm] (sembed1) at (35,13) {};
\filldraw[draw=NNInput,fill=NNInput!18,line width=.5pt] (29.0,11.5) rectangle ++(3.0,3);
\filldraw[draw=NNInput,fill=NNInput!18,line width=.5pt] (32.0,11.5) rectangle ++(3.0,3);
\filldraw[draw=NNInput,fill=NNInput!18,line width=.5pt] (35.0,11.5) rectangle ++(3.0,3);
\filldraw[draw=NNInput,fill=NNInput!18,line width=.5pt] (38.0,11.5) rectangle ++(3.0,3);
\draw[wire] (stoken1.north)--(sembed1.south);
\draw[wire] (sembed1.north)--(35,51);
\node[] (stoken2) at (53,0) {狗};
\node[inner sep=0pt,minimum width=12mm,minimum height=3mm] (sembed2) at (53,13) {};
\filldraw[draw=NNInput,fill=NNInput!18,line width=.5pt] (47.0,11.5) rectangle ++(3.0,3);
\filldraw[draw=NNInput,fill=NNInput!18,line width=.5pt] (50.0,11.5) rectangle ++(3.0,3);
\filldraw[draw=NNInput,fill=NNInput!18,line width=.5pt] (53.0,11.5) rectangle ++(3.0,3);
\filldraw[draw=NNInput,fill=NNInput!18,line width=.5pt] (56.0,11.5) rectangle ++(3.0,3);
\draw[wire] (stoken2.north)--(sembed2.south);
\draw[wire] (sembed2.north)--(53,51);
\node[stage,minimum width=58mm,minimum height=10mm] (slayer1) at (35,57) {第$1$层：双向注意力＋FFN};
\node[stage,minimum width=58mm,minimum height=10mm] (slayerL) at (35,89) {第$L$层：双向注意力＋FFN};
\draw[wire] (17,62)--(17,67);
\node[] (sdots17) at (17,73) {$\vdots$};
\draw[wire] (17,79)--(17,84);
\draw[wire] (35,62)--(35,67);
\node[] (sdots35) at (35,73) {$\vdots$};
\draw[wire] (35,79)--(35,84);
\draw[wire] (53,62)--(53,67);
\node[] (sdots53) at (53,73) {$\vdots$};
\draw[wire] (53,79)--(53,84);
\node[nn annotation] (sbias) at (35,132) {注意力内加入相对位置偏置};
\node[] (ttoken0) at (106,0) {START};
\node[inner sep=0pt,minimum width=12mm,minimum height=3mm] (tembed0) at (106,13) {};
\filldraw[draw=NNInput,fill=NNInput!18,line width=.5pt] (100.0,11.5) rectangle ++(3.0,3);
\filldraw[draw=NNInput,fill=NNInput!18,line width=.5pt] (103.0,11.5) rectangle ++(3.0,3);
\filldraw[draw=NNInput,fill=NNInput!18,line width=.5pt] (106.0,11.5) rectangle ++(3.0,3);
\filldraw[draw=NNInput,fill=NNInput!18,line width=.5pt] (109.0,11.5) rectangle ++(3.0,3);
\draw[wire] (ttoken0.north)--(tembed0.south);
\draw[wire] (tembed0.north)--(106,51);
\node[] (ttoken1) at (129,0) {$\langle S_1\rangle$};
\node[inner sep=0pt,minimum width=12mm,minimum height=3mm] (tembed1) at (129,13) {};
\filldraw[draw=NNInput,fill=NNInput!18,line width=.5pt] (123.0,11.5) rectangle ++(3.0,3);
\filldraw[draw=NNInput,fill=NNInput!18,line width=.5pt] (126.0,11.5) rectangle ++(3.0,3);
\filldraw[draw=NNInput,fill=NNInput!18,line width=.5pt] (129.0,11.5) rectangle ++(3.0,3);
\filldraw[draw=NNInput,fill=NNInput!18,line width=.5pt] (132.0,11.5) rectangle ++(3.0,3);
\draw[wire] (ttoken1.north)--(tembed1.south);
\draw[wire] (tembed1.north)--(129,51);
\node[] (ttoken2) at (152,0) {追};
\node[inner sep=0pt,minimum width=12mm,minimum height=3mm] (tembed2) at (152,13) {};
\filldraw[draw=NNInput,fill=NNInput!18,line width=.5pt] (146.0,11.5) rectangle ++(3.0,3);
\filldraw[draw=NNInput,fill=NNInput!18,line width=.5pt] (149.0,11.5) rectangle ++(3.0,3);
\filldraw[draw=NNInput,fill=NNInput!18,line width=.5pt] (152.0,11.5) rectangle ++(3.0,3);
\filldraw[draw=NNInput,fill=NNInput!18,line width=.5pt] (155.0,11.5) rectangle ++(3.0,3);
\draw[wire] (ttoken2.north)--(tembed2.south);
\draw[wire] (tembed2.north)--(152,51);
\node[stage,minimum width=66mm,minimum height=10mm] (tlayer1) at (129,57) {第$1$层：因果＋交叉注意力＋FFN};
\node[stage,minimum width=66mm,minimum height=10mm] (tlayerL) at (129,89) {第$L$层：因果＋交叉注意力＋FFN};
\draw[wire] (106,62)--(106,67);
\node[] (tdots106) at (106,73) {$\vdots$};
\draw[wire] (106,79)--(106,84);
\draw[wire] (129,62)--(129,67);
\node[] (tdots129) at (129,73) {$\vdots$};
\draw[wire] (129,79)--(129,84);
\draw[wire] (152,62)--(152,67);
\node[] (tdots152) at (152,73) {$\vdots$};
\draw[wire] (152,79)--(152,84);
\node[nn annotation] (tbias) at (129,132) {注意力内加入相对位置偏置};
\draw[wire] (17,94)--(17,107);
\node[inner sep=0pt,minimum width=12mm,minimum height=3mm] (sh0) at (17,109) {};
\filldraw[draw=NNOutput,fill=NNOutput!18,line width=.5pt] (11.0,107.5) rectangle ++(3.0,3);
\filldraw[draw=NNOutput,fill=NNOutput!18,line width=.5pt] (14.0,107.5) rectangle ++(3.0,3);
\filldraw[draw=NNOutput,fill=NNOutput!18,line width=.5pt] (17.0,107.5) rectangle ++(3.0,3);
\filldraw[draw=NNOutput,fill=NNOutput!18,line width=.5pt] (20.0,107.5) rectangle ++(3.0,3);
\draw[wire] (35,94)--(35,107);
\node[inner sep=0pt,minimum width=12mm,minimum height=3mm] (sh1) at (35,109) {};
\filldraw[draw=NNOutput,fill=NNOutput!18,line width=.5pt] (29.0,107.5) rectangle ++(3.0,3);
\filldraw[draw=NNOutput,fill=NNOutput!18,line width=.5pt] (32.0,107.5) rectangle ++(3.0,3);
\filldraw[draw=NNOutput,fill=NNOutput!18,line width=.5pt] (35.0,107.5) rectangle ++(3.0,3);
\filldraw[draw=NNOutput,fill=NNOutput!18,line width=.5pt] (38.0,107.5) rectangle ++(3.0,3);
\draw[wire] (53,94)--(53,107);
\node[inner sep=0pt,minimum width=12mm,minimum height=3mm] (sh2) at (53,109) {};
\filldraw[draw=NNOutput,fill=NNOutput!18,line width=.5pt] (47.0,107.5) rectangle ++(3.0,3);
\filldraw[draw=NNOutput,fill=NNOutput!18,line width=.5pt] (50.0,107.5) rectangle ++(3.0,3);
\filldraw[draw=NNOutput,fill=NNOutput!18,line width=.5pt] (53.0,107.5) rectangle ++(3.0,3);
\filldraw[draw=NNOutput,fill=NNOutput!18,line width=.5pt] (56.0,107.5) rectangle ++(3.0,3);
\node[nn heading] (hs) at (35,121) {编码器最终输出 $\bm H_{\mathrm s}$};
\draw[draw=NNWire,line width=1.2pt] (17,115)--(79,115)--(79,57);
\draw[draw=NNWire,line width=1.2pt] (17,110.5)--(17,115);
\draw[draw=NNWire,line width=1.2pt] (35,110.5)--(35,115);
\draw[draw=NNWire,line width=1.2pt] (53,110.5)--(53,115);
\draw[wire] (79,57)--(tlayer1.west);
\fill[NNWire] (79,57) circle (.55);
\node[nn annotation] (kv1) at (87,63) {K、V};
\draw[wire] (79,89)--(tlayerL.west);
\fill[NNWire] (79,89) circle (.55);
\node[nn annotation] (kvL) at (87,95) {K、V};
\node[stage,draw=NNOutput,fill=NNOutput!18,minimum width=14mm,minimum height=8mm] (lm0) at (106,108) {LM};
\draw[wire] (106,94)--(lm0.south);
\node[] (prediction0) at (106,124) {$\langle S_1\rangle$};
\draw[wire] (lm0.north)--(prediction0.south);
\node[stage,draw=NNOutput,fill=NNOutput!18,minimum width=14mm,minimum height=8mm] (lm1) at (129,108) {LM};
\draw[wire] (129,94)--(lm1.south);
\node[] (prediction1) at (129,124) {追};
\draw[wire] (lm1.north)--(prediction1.south);
\node[stage,draw=NNOutput,fill=NNOutput!18,minimum width=14mm,minimum height=8mm] (lm2) at (152,108) {LM};
\draw[wire] (152,94)--(lm2.south);
\node[] (prediction2) at (152,124) {$\langle S_2\rangle$};
\draw[wire] (lm2.north)--(prediction2.south);
\node[nn heading] (enc) at (35,141) {编码器：读取全部源词元};
\node[nn heading] (dec) at (129,141) {解码器：读取已知目标前缀};
\node[nn annotation] (inputnote) at (35,-13) {源输入};
\node[nn annotation] (targetnote) at (129,-13) {右移后的目标输入};
\node[nn annotation] (knot) at (129,155) {交叉注意力的Q来自目标路径};
\end{tikzpicture}
